\documentclass[12pt,a4paper]{amsart}

\usepackage[ngerman]{babel}
\usepackage{amsmath,amssymb,amsthm}
\usepackage{mathtools}
\usepackage[colorlinks=true,linkcolor=blue,citecolor=blue]{hyperref}
\usepackage{geometry}
\geometry{margin=2.5cm}

% -------------------------------------------------------
% Theorem-Umgebungen
% -------------------------------------------------------
\newtheorem{theorem}{Satz}[section]
\newtheorem{lemma}[theorem]{Lemma}
\newtheorem{korollar}[theorem]{Korollar}
\newtheorem{proposition}[theorem]{Proposition}
\theoremstyle{definition}
\newtheorem{definition}[theorem]{Definition}
\newtheorem{bemerkung}[theorem]{Bemerkung}

% -------------------------------------------------------
% Eigene Makros
% -------------------------------------------------------
\newcommand{\N}{\mathbb{N}}
\newcommand{\Z}{\mathbb{Z}}
\newcommand{\Q}{\mathbb{Q}}
\newcommand{\Fp}{\mathbb{F}_p}
\DeclareMathOperator{\ord}{ord}
\DeclareMathOperator{\ggT}{ggT}
\DeclareMathOperator{\kgV}{kgV}

% -------------------------------------------------------
\begin{document}
% -------------------------------------------------------

\title{Der Satz von Wilson für Primzahlpotenzen}

\author{Michael Fuhrmann}
\address{Specialist-Mathematics-Projekt, Build 63}
\email{specialist-maths@localhost}
\date{\today}

\begin{abstract}
Wir untersuchen das Produkt aller Einheiten im Ring $\Z/p^k\Z$ und beweisen
das folgende vollständige Ergebnis: Für eine ungerade Primzahl $p$ und jedes
$k \geq 1$ gilt
\[
  \prod_{\substack{a=1\\\ggT(a,p^k)=1}}^{p^k} a \;\equiv\; -1 \pmod{p^k}.
\]
Für $p = 2$ ist die Situation besonders: Das analoge Produkt ist für $k \geq 3$
gleich $1 \pmod{2^k}$, während die kleinen Fälle $k=1$ und $k=2$ separat
behandelt werden.
Der Beweis beruht auf zwei Säulen: (i) der Paarungsmethode – jede Einheit wird
mit ihrem Inversen gepaart – und (ii) der Tatsache, dass die einzigen
Quadratwurzeln der Eins modulo $p^k$ (für ungerades $p$) genau $\pm 1$ sind.
Als Korollar identifizieren wir die Struktur der Einheitengruppe $(\Z/p^k\Z)^*$.
\end{abstract}

\maketitle
\tableofcontents

% ================================================================
\section{Einleitung}
% ================================================================

Der klassische Satz von Wilson besagt, dass $(p-1)! \equiv -1 \pmod{p}$ für
jede Primzahl $p$ gilt. Das Produkt $(p-1)!$ ist nichts anderes als das
Produkt \emph{aller} Elemente der multiplikativen Gruppe $(\Z/p\Z)^*$.
Eine natürliche Frage lautet: Was passiert, wenn wir den Modul $p$ durch eine
Primzahlpotenz $p^k$ ersetzen?

Die Einheitengruppe $(\Z/p^k\Z)^*$ besteht aus allen Restklassen $1 \leq a \leq p^k$
mit $\ggT(a, p^k) = 1$ und hat Ordnung $\varphi(p^k) = p^{k-1}(p-1)$.
Für ungerade Primzahlen $p$ ist diese Gruppe zyklisch, und das Produkt aller
ihrer Elemente lässt sich allein aus der Gruppenstruktur berechnen.
Für $p = 2$ ist die Gruppe für $k \geq 3$ nicht mehr zyklisch, und das
Produkt ist $+1$ statt $-1$.

Im gesamten Artikel bezeichnet $p$ eine Primzahl und $k$ eine positive ganze
Zahl. Alle Kongruenzen gelten modulo $p^k$, sofern nicht anders angegeben.

% ================================================================
\section{Quadratwurzeln der Eins modulo \texorpdfstring{$p^k$}{p hoch k}}
% ================================================================

Das folgende Schlüssellemma identifiziert die selbst-inversen Elemente von
$(\Z/p^k\Z)^*$.

\begin{lemma}[Quadratwurzeln der Eins für ungerade Primzahlpotenzen]%
  \label{lem:quadratwurzel_eins}
  Sei $p$ eine ungerade Primzahl und $k \geq 1$. Die Lösungen von
  $a^2 \equiv 1 \pmod{p^k}$ in $\{1, \ldots, p^k\}$ sind genau
  $a = 1$ und $a = p^k - 1 \equiv -1 \pmod{p^k}$.
\end{lemma}

\begin{proof}
  Es gilt $a^2 \equiv 1 \pmod{p^k}$ genau dann, wenn $p^k \mid (a-1)(a+1)$.
  Da $\ggT(a-1, a+1) \mid 2$ und $p$ ungerade ist, kann $p$ nicht gleichzeitig
  $a-1$ und $a+1$ teilen. Daher teilt $p^k$ entweder $a-1$ oder $a+1$, was
  $a \equiv 1 \pmod{p^k}$ oder $a \equiv -1 \pmod{p^k}$ ergibt.

  Genauer: Angenommen $p^k \mid (a-1)(a+1)$. Schreibe $a - 1 = p^\alpha u$
  und $a + 1 = p^\beta v$ mit $\ggT(u,p) = \ggT(v,p) = 1$.
  Dann gilt $(a+1) - (a-1) = 2$, also $p^{\min(\alpha,\beta)} \mid 2$.
  Da $p$ ungerade ist, folgt $\min(\alpha, \beta) = 0$.
  Daher gilt entweder $\alpha = 0$ (d.\,h.\ $p \nmid a-1$) und $p^k \mid a+1$,
  oder $\beta = 0$ (d.\,h.\ $p \nmid a+1$) und $p^k \mid a-1$.
  In beiden Fällen sind die einzigen Lösungen $a \equiv \pm 1 \pmod{p^k}$.
\end{proof}

\begin{bemerkung}
  Lemma~\ref{lem:quadratwurzel_eins} gilt nicht für $p = 2$ und $k \geq 3$:
  Die Kongruenz $a^2 \equiv 1 \pmod{8}$ hat vier Lösungen modulo $8$:
  $a \in \{1, 3, 5, 7\}$. Allgemein hat $a^2 \equiv 1 \pmod{2^k}$
  für $k \geq 3$ genau vier Lösungen: $a \equiv \pm 1 \pmod{2^{k-1}}$.
\end{bemerkung}

% ================================================================
\section{Hauptsatz für ungerade Primzahlen}
% ================================================================

\begin{theorem}[Wilson für Primzahlpotenzen, ungerades $p$]\label{satz:wilson_pk_ungerade}
  Sei $p$ eine ungerade Primzahl und $k \geq 1$. Dann gilt
  \[
    \prod_{\substack{a=1\\\ggT(a,p^k)=1}}^{p^k} a
    \;\equiv\; -1 \pmod{p^k}.
  \]
\end{theorem}

\begin{proof}
  Bezeichne mit $P$ das Produkt auf der linken Seite. Wir arbeiten in der
  Gruppe $G = (\Z/p^k\Z)^*$ der Ordnung $\varphi(p^k) = p^{k-1}(p-1)$.

  \textbf{Schritt 1 (Paarung).}
  Für jedes $a \in G$ bilden wir das Paar $\{a, a^{-1}\}$ mit dem
  multiplikativen Inversen $a^{-1} \pmod{p^k}$. Falls $a \neq a^{-1}$,
  ist $\{a, a^{-1}\}$ eine zweielementige Teilmenge von $G$ und trägt den
  Faktor $a \cdot a^{-1} = 1$ zu $P$ bei.

  \textbf{Schritt 2 (Selbst-inverse Elemente).}
  Ein Element $a$ ist selbst-invers (d.\,h.\ $a = a^{-1}$) genau dann, wenn
  $a^2 \equiv 1 \pmod{p^k}$. Nach Lemma~\ref{lem:quadratwurzel_eins} sind die
  einzigen selbst-inversen Elemente $a \equiv 1$ und $a \equiv -1 \pmod{p^k}$.

  \textbf{Schritt 3 (Berechnung des Produkts).}
  Nach der Paarung bilden alle Elemente von $G \setminus \{1, -1\}$
  insgesamt $(|G| - 2)/2 = (p^{k-1}(p-1) - 2)/2$ Paare, von denen jedes
  den Faktor $1$ beiträgt. Die verbleibenden Faktoren stammen von den beiden
  selbst-inversen Elementen:
  \[
    P \;\equiv\; 1 \cdot (-1) \;\equiv\; -1 \pmod{p^k}.
  \]

  \textbf{Schritt 4 (Gültigkeit der Paarungsanzahl).}
  Wir müssen überprüfen, dass $|G| - 2$ gerade ist, d.\,h.\ dass
  $|G| = p^{k-1}(p-1)$ gerade ist. Da $p$ eine ungerade Primzahl ist,
  ist $p - 1$ gerade, also ist $p^{k-1}(p-1)$ tatsächlich gerade.
  Die Paarung ist daher korrekt und die obige Rechnung gültig.
\end{proof}

% ================================================================
\section{Der Fall \texorpdfstring{$p = 2$}{p = 2}}
% ================================================================

Die Primzahl $p = 2$ erfordert eine gesonderte Analyse, da die Gruppe
$(\Z/2^k\Z)^*$ für $k \geq 3$ nicht zyklisch ist.

\begin{theorem}[Wilson für Zweierpotenzen]\label{satz:wilson_pk_2}
  Sei $k \geq 1$. Dann gilt
  \[
    \prod_{\substack{a=1\\\ggT(a,2^k)=1}}^{2^k} a
    \;\equiv\;
    \begin{cases}
      1  \pmod{2}   & k = 1, \\
      -1 \pmod{4}   & k = 2, \\
      1  \pmod{2^k} & k \geq 3.
    \end{cases}
  \]
\end{theorem}

\begin{proof}
  \textbf{Fall $k = 1$.}
  Die einzige ungerade Restklasse modulo $2$ in $\{1\}$ ist $1$.
  Das Produkt ist $1 \equiv 1 \pmod{2}$.

  \textbf{Fall $k = 2$ (modulo $4$).}
  Die Einheiten modulo $4$ sind $\{1, 3\}$.
  Ihr Produkt ist $1 \cdot 3 = 3 \equiv -1 \pmod{4}$.

  \textbf{Fall $k \geq 3$.}
  Die Einheitengruppe $(\Z/2^k\Z)^*$ hat Ordnung $\varphi(2^k) = 2^{k-1}$.
  Für $k \geq 3$ ist sie \emph{nicht} zyklisch; stattdessen gilt
  \[
    (\Z/2^k\Z)^* \;\cong\; \Z/2 \times \Z/2^{k-2}.
  \]
  Die Elemente der Ordnung höchstens $2$ (d.\,h.\ Lösungen von
  $a^2 \equiv 1 \pmod{2^k}$) sind $a \equiv \pm 1 \pmod{2^{k-1}}$,
  was vier Lösungen modulo $2^k$ ergibt:
  $1$, $2^{k-1}-1$, $2^{k-1}+1$ und $2^k - 1 \equiv -1$.

  Die Untergruppe der selbst-inversen Elemente ist also
  $S = \{1,\; -1,\; 2^{k-1}-1,\; 2^{k-1}+1\} \pmod{2^k}$
  mit Ordnung $4$. Das Produkt aller Elemente von $S$ ist:
  \[
    1 \cdot (-1) \cdot (2^{k-1}-1) \cdot (2^{k-1}+1)
    = (-1)\bigl((2^{k-1})^2 - 1\bigr)
    = -(2^{2k-2} - 1).
  \]
  Wir berechnen dies modulo $2^k$: Da $k \geq 3$, gilt $2k-2 \geq k+1 > k$,
  also $2^{2k-2} \equiv 0 \pmod{2^k}$, und
  \[
    -(2^{2k-2}-1) \equiv -(-1) = 1 \pmod{2^k}.
  \]
  Alle nicht-selbst-inversen Elemente bilden Paare $\{a, a^{-1}\}$, die je
  den Faktor $1$ beitragen. Das Gesamtprodukt ist daher $P \equiv 1 \pmod{2^k}$.
\end{proof}

% ================================================================
\section{Struktur der Einheitengruppe und Korollare}
% ================================================================

Der Beweis von Satz~\ref{satz:wilson_pk_ungerade} beruht auf der Zyklizität
von $(\Z/p^k\Z)^*$ für ungerades $p$, die wir nun präzise formulieren.

\begin{theorem}[Struktur von $(\Z/p^k\Z)^*$]\label{satz:struktur}
  \begin{enumerate}
    \item[(a)] Für jede ungerade Primzahl $p$ und $k \geq 1$:
      $(\Z/p^k\Z)^* \cong \Z_{\varphi(p^k)} = \Z_{p^{k-1}(p-1)}$
      (zyklisch der Ordnung $p^{k-1}(p-1)$).
    \item[(b)] Für $p = 2$: $(\Z/2\Z)^* = \{1\}$, $(\Z/4\Z)^* \cong \Z_2$,
      und $(\Z/2^k\Z)^* \cong \Z_2 \times \Z_{2^{k-2}}$ für $k \geq 3$.
  \end{enumerate}
\end{theorem}

\begin{proof}
  \textbf{Teil (a).}
  Wir zeigen, dass $(\Z/p^k\Z)^*$ zyklisch ist, indem wir eine Primitivwurzel
  modulo $p$ zu einer Primitivwurzel modulo $p^k$ liften.

  Sei $g$ eine Primitivwurzel modulo $p$ (die existiert, da $(\Z/p\Z)^*$
  zyklisch von Ordnung $p-1$ ist). Wir können $g$ so wählen, dass
  $g^{p-1} \not\equiv 1 \pmod{p^2}$ gilt. (Falls dies für $g$ nicht gilt,
  ersetze $g$ durch $g + p$, das noch $(\Z/p\Z)^*$ erzeugt; höchstens ein
  Wert kann versagen, da es $p-1$ Kandidaten gibt.)

  Wir behaupten, dass $g$ für alle $k \geq 1$ eine Primitivwurzel modulo
  $p^k$ ist. Durch Induktion: Angenommen, $g$ hat die Ordnung
  $p^{k-1}(p-1)$ modulo $p^k$. Schreibe
  \[
    g^{p-1} = 1 + p \cdot t, \quad t \not\equiv 0 \pmod{p}.
  \]
  Durch den Binomialsatz (oder Hensel-Lifting) ergibt sich
  $g^{p^{k-1}(p-1)} = (1+pt)^{p^{k-1}} \equiv 1 + p^k t' \pmod{p^{k+1}}$
  mit $t' \not\equiv 0 \pmod{p}$.
  Also hat $g$ modulo $p^{k+1}$ die Ordnung $p^k(p-1) = \varphi(p^{k+1})$,
  was beweist, dass $g$ auch eine Primitivwurzel modulo $p^{k+1}$ ist.

  \textbf{Teil (b).}
  Die Einheitengruppe $(\Z/2\Z)^* = \{1\}$ ist trivial.
  $(\Z/4\Z)^* = \{1, 3\} \cong \Z_2$.
  Für $k \geq 3$: Das Element $-1 \equiv 2^k - 1$ erzeugt eine Untergruppe
  der Ordnung $2$, und $5$ erzeugt eine Untergruppe der Ordnung $2^{k-2}$
  (da $5^{2^{k-3}} \equiv 1 + 2^{k-1} \not\equiv 1 \pmod{2^k}$, aber
  $5^{2^{k-2}} \equiv 1 \pmod{2^k}$). Zusammen erzeugen $-1$ und $5$ die
  gesamte Gruppe, und $\langle -1 \rangle \cap \langle 5 \rangle = \{1\}$,
  was die direkte Produktzerlegung $(\Z/2^k\Z)^* \cong \Z_2 \times \Z_{2^{k-2}}$
  liefert.
\end{proof}

\begin{korollar}[Eindeutiges Element der Ordnung 2 für ungerade Primzahlpotenzen]%
  \label{kor:eindeutig_ordnung2}
  Für eine ungerade Primzahl $p$ und $k \geq 1$ hat die Gruppe $(\Z/p^k\Z)^*$
  genau ein Element der Ordnung $2$, nämlich $-1 \equiv p^k - 1 \pmod{p^k}$.
\end{korollar}

\begin{proof}
  Nach Satz~\ref{satz:struktur}(a) ist $(\Z/p^k\Z)^*$ zyklisch von gerader
  Ordnung $p^{k-1}(p-1)$. In jeder zyklischen Gruppe gerader Ordnung gibt es
  genau ein Element der Ordnung $2$, nämlich das eindeutige Element $\tau$
  mit $\tau^2 = e$ und $\tau \neq e$. Nach Lemma~\ref{lem:quadratwurzel_eins}
  ist dieses Element $-1$.
\end{proof}

\begin{korollar}[Wilson-Produktformel]\label{kor:wilson_formel}
  Für eine ungerade Primzahl $p$ und $k \geq 1$ gilt:
  \[
    \prod_{\substack{a=1\\\ggT(a,p^k)=1}}^{p^k} a \;\equiv\; -1 \pmod{p^k}.
  \]
  Für $k = 1$ ergibt sich der klassische Satz von Wilson: $(p-1)! \equiv -1 \pmod{p}$.
\end{korollar}

\begin{proof}
  Dies ist Satz~\ref{satz:wilson_pk_ungerade}. Für $k = 1$ ist das Produkt
  über $\{a : 1 \leq a \leq p,\ \ggT(a,p) = 1\}$ gleich $(p-1)!$.
\end{proof}

% ================================================================
\section{Schlussfolgerung}
% ================================================================

Wir haben eine vollständige Verallgemeinerung des Satzes von Wilson auf
Primzahlpotenzen bewiesen. Die entscheidende Erkenntnis ist, dass für ungerade
Primzahlpotenzen $p^k$ die Einheitengruppe zyklisch ist und genau \emph{ein}
selbst-inverses Element außer dem Neutralelement besitzt, nämlich $-1$.
Dies zwingt das Produkt aller Einheiten, gleich $-1$ zu sein.

Für $p = 2$ führt das Versagen der Zyklizität für $k \geq 3$ zu vier
selbst-inversen Elementen, deren gemeinsames Produkt $+1$ ergibt, was
das Vorzeichen im Vergleich zum ungerade-Prim-Fall umkehrt.

\begin{thebibliography}{9}

\bibitem{Hardy1979}
G.~H.~Hardy und E.~M.~Wright.
\newblock \emph{Einführung in die Zahlentheorie} (An Introduction to the Theory of Numbers), 5.~Aufl.
\newblock Oxford University Press, 1979.

\bibitem{Ireland1990}
K.~Ireland und M.~Rosen.
\newblock \emph{A Classical Introduction to Modern Number Theory}, 2.~Aufl.
\newblock Springer, New York, 1990.

\bibitem{Niven1991}
I.~Niven, H.~S.~Zuckerman und H.~L.~Montgomery.
\newblock \emph{An Introduction to the Theory of Numbers}, 5.~Aufl.
\newblock Wiley, New York, 1991.

\bibitem{Rose1994}
H.~E.~Rose.
\newblock \emph{A Course in Number Theory}, 2.~Aufl.
\newblock Oxford University Press, 1994.

\end{thebibliography}

\end{document}
